\scanpage{37}
\mspage{6}
\noindent (19)
\diagram{\begin{tikzcd}[row sep=large]
\mathbb D^*_a \arrow[r,hook,"{\text{``bouchage du trou }a\text{''}}"] \arrow[d] & \mathbb D_a \\
X^*=X\setminus\{a\}
\end{tikzcd}}

Ici le diagramme (10) se réduit à

\noindent (20)
\diagram{\begin{tikzcd}
\mathbb Z_{a,X} \arrow[r] \arrow[d,"{\deleted{\unsure{\varphi}}}"] & \{1\} \\
\pi_1(X^*)
\end{tikzcd}},

où $\mathbb Z_{a,X}$\inserted{$\simeq\mathbb Z$} est le ``module d'orientation de
$X$ en $a$'' dont les deux gén. correspondent
aux deux orientations de $X$ en $a$. L'hom.
vertical de \unsure{(2)|(20)} est l'hom. qui définit
la structure ``lacets en $a$'' du groupe
fondamental $X^*$. Le cas le plus
simple est celui où
\[
(21)\qquad\left\{\begin{array}{l}
X=\mathbb C\quad\text{plan complexe}\\
a=\{0\}
\end{array}\right.
\]
\margin{d'où \diagram{\begin{tikzcd}[column sep=small,row sep=small]
\mathbb D^*_0\arrow[r]\arrow[d]&\mathbb D_0\\
\mathbb C^*
\end{tikzcd}}}
dont le produit avec lui-même un nb
arbitraire de fois, donne les ``modèles-
types'' des stratifications associées aux
diviseurs à croisements normaux (dans le
\mspage{36}
